\documentclass[11pt,a4paper]{article}
\usepackage[utf8]{inputenc}
\usepackage[T1]{fontenc}
\usepackage[ngerman]{babel}
\usepackage{lmodern,amsmath,amssymb,amsthm,booktabs}
\usepackage[margin=27mm]{geometry}
\usepackage[hidelinks]{hyperref}
\newtheorem{satz}{Satz}
\newtheorem{lemma}{Lemma}
\title{Eine Randfall-Widerlegung der verallgemeinerten\newline Baloui-Vermutung in OEIS A381054}
\author{Arbeitsnotiz für Jonas J. Gebendorfer}
\date{13. Oktober 2026}
\begin{document}
\maketitle
\noindent\textbf{Ergebnis.} Die in den Kommentaren von A381054 formulierte allgemeine
Äquivalenz ist falsch. Für den dort zulässigen Primzahlfall $m=3$, $n=11$ gilt
\[
 a_3(11)=37\ne 0=(\lfloor\sqrt{11}\rfloor-3)^2.
\]
Der Abschluss betrifft genau diese veröffentlichte Fassung. Die Ursache ist eine
unzureichende untere Primzahlgrenze. Der folgende Handbeweis benötigt keine
numerische Näherung, keine endliche Suchraumannahme und keine externe mathematische
Prämisse. Neuheit und externes Review sind nicht festgestellt.

\section{Originalziel und Definitionen}
Die am 13.10.2026 abgerufene Originalseite \href{https://oeis.org/A381054}{OEIS A381054}
von Hassan Baloui definiert die Familie über den jeweils kleinsten positiven
Multiplikator. Setze für $q\ge1$
\[
 d(q)=\max\{d\in\mathbb Z_{\ge1}:d\mid q,\ d^2\le q\},\qquad
 C(q)=\left\lfloor\sqrt{q/d(q)}\right\rfloor-\left\lfloor\sqrt{d(q)}\right\rfloor,
\]
und, wenn die Menge nicht leer ist,
\[
 a_m(n)=\min\{k\in\mathbb Z_{\ge1}:C(nk)=m\}.
\]
Der Index A000196 bezeichnet ausdrücklich $\lfloor\sqrt n\rfloor$.
Die verallgemeinerte Vermutung im Kommentar verlangt für positive ganze $m,n$
\begin{equation}\label{original}
 a_m(n)=(\lfloor\sqrt n\rfloor-m)^2
 \quad\Longleftrightarrow\quad
 n\text{ ist prim und }n>m^2+1.
\end{equation}
Die Seite nennt die Fälle $m=0,1$ bereits beantwortet. Das heutige Ziel ist die
allgemeine Aussage, konkret ihre Richtung von der Primzahlbedingung zur Gleichheit
im Fall $m=3$. Die Programme der Quelle beginnen die Multiplikatorsuche bei $k=1$.
Eine Lesart mit $k=0$ würde zudem $d(0)$ verlangen und entspricht dieser Definition nicht.

\section{Vollständiger Gegenbeweis}
\begin{satz}\label{counter}
Es gilt $a_3(11)=37$. Damit ist die universelle Aussage \eqref{original} widerlegt.
\end{satz}
\begin{proof}
Die Zahl $11$ ist prim und $11>3^2+1=10$. Da
$3^2\le11<4^2$, ist der von \eqref{original} verlangte Wert
$(3-3)^2=0$.

Wir bestimmen das tatsächlich definierte Minimum.
Für $1\le k<11$ gilt $d(11k)=k$: Einerseits ist $k<\sqrt{11k}<11$;
andererseits ist jeder Teiler von $11k$ unterhalb von $11$ wegen der Primheit von
$11$ ein Teiler von $k$. Daher
\[
 C(11k)=3-\lfloor\sqrt k\rfloor\le2 \qquad(1\le k\le10).
\]
Für $11\le k\le35$ ist $11\le\sqrt{11k}$, also $d(11k)\ge11$.
Hieraus folgen $11k/d(11k)\le k\le35$ und
\[
 C(11k)\le\lfloor\sqrt{35}\rfloor-\lfloor\sqrt{11}\rfloor
 =5-3=2.
\]
Für $k=36$ ist $11k=396$. Die Teiler bis $\lfloor\sqrt{396}\rfloor=19$
sind genau $1,2,3,4,6,9,11,12,18$. Somit $d(396)=18$, $396/d(396)=22$ und
\[
 C(396)=\lfloor\sqrt{22}\rfloor-\lfloor\sqrt{18}\rfloor=4-4=0.
\]
Für $k=37$ ist $407=11\cdot37$. Die Faktoren $11$ und $37$ sind prim,
$11^2<407<37^2$, und daher $d(407)=11$. Folglich
\[
 C(407)=\lfloor\sqrt{37}\rfloor-\lfloor\sqrt{11}\rfloor=6-3=3.
\]
Es existiert also ein zulässiger Multiplikator, und keiner der kleineren positiven
Multiplikatoren ist zulässig. Damit $a_3(11)=37$.
Die rechte Seite von \eqref{original} ist wahr und die linke Seite falsch.
Ein einziger solcher Fall negiert die gesamte universelle Äquivalenz.
\end{proof}

Der noch kürzere Widerspruch benötigt nur die Existenz des Multiplikators $37$:
Ein Minimum über positive ganze Zahlen kann niemals $0$ sein. Die Minimalitätsprüfung
oben erhält zusätzlich den exakten Wert und schützt die Aufgabenabbildung.

\section{Was an der Primzahlformel trotzdem trägt}
\begin{lemma}\label{prime}
Sei $p$ prim, $s=\lfloor\sqrt p\rfloor$ und $1\le m<s$.
Dann ist $a_m(p)$ definiert und
\[
 a_m(p)=(s-m)^2.
\]
\end{lemma}
\begin{proof}
Setze $k_0=(s-m)^2$. Wegen $1\le m<s$ gilt $1\le k_0<p$.
Für jedes $1\le k\le k_0$ ist $d(pk)=k$: Ein Teiler unter
$\sqrt{pk}<p$ kann den Primfaktor $p$ nicht enthalten, und $k$ selbst liegt unter
$\sqrt{pk}$. Also $C(pk)=s-\lfloor\sqrt k\rfloor$.
Bei $k=k_0$ ist dies $m$. Bei $k<k_0$ gilt
$\lfloor\sqrt k\rfloor\le s-m-1$, also $C(pk)\ge m+1$.
Das beweist das Minimum.
\end{proof}
Die notwendige Bereichsbedingung für diese Herleitung ist $m<s$, äquivalent zu
$p\ge(m+1)^2$. Die in der Quelle verwendete Bedingung $p>m^2+1$ reicht dafür
nicht aus. Die Lemma-Aussage belegt eine gültige Primzahlrichtung im engeren Bereich;
sie entscheidet keinen allgemeinen Umkehrschluss für zusammengesetzte Zahlen.
Sie wird nicht als neuer vollständiger Ersatz der Quellvermutung ausgegeben.

\section{Quellenstand, Kontrollen und Reichweite}
Die Original- und Internal-Seite von A381054 wurden direkt per HTTP abgerufen.
Die Internal-Fassung trägt Revision \texttt{\#43}, 5. Juni 2025.
Der am 13.10.2026 erhaltene HTML-Snapshot der Originalseite hat SHA-256
\begin{center}\small\ttfamily
29ef6fff68532a0cb7d46953dea656a30616c1c7c5f258a58e7484a20ee57eb0
\end{center}
Die passende Quellstelle ist der Familienkommentar unmittelbar vor den Links;
Definition und Startwert werden durch die Maple-/PARI-Abschnitte bestätigt.
Gezielte Suchwege fanden keine früher gespeicherte eigene Behandlung dieser Familie
und keinen anderen Schluss des hier geprüften allgemeinen Kommentars. Das ist
kein weltweiter Prioritätsnachweis.

\texttt{verify.py} verwendet Teilerprüfung bis zur ganzzahligen Quadratwurzel.
\texttt{verify.js} faktorisiert die Zahlen, erzeugt alle Teiler und nimmt den mittleren
Teiler; die Quadratwurzeln werden ganzzahlig per Binärsuche ermittelt.
Beide Kontrollen geben für alle $1\le k\le37$ dieselben Paare und Lücken aus.
Positive Kontrollen, $k=0$, $k=36$ und ein absichtlich nicht nächstes Faktorenpaar
werden geprüft. Insbesondere würde das falsche Paar $(11,36)$ für $396$ die Lücke
$3$ vortäuschen; das korrekte Paar $(18,22)$ hat Lücke $0$.
Diese Prüfungen sind intern und teilen Quelle und Bearbeiter. Sie werden nicht
als externes unabhängiges Review ausgegeben. Der universelle Negationsschluss trägt
durch Satz~\ref{counter}, nicht durch die Wiederholung eines Programms.

\medskip
\noindent\textbf{Zielwahl.} A383115 war ein zuvor erwogener Kandidat. Vor dem Angriff
wurde ein öffentlicher Beweisvorschlag gefunden
(\href{https://math.stackexchange.com/questions/5148643/}{Math.SE 5148643}).
Diese bereits veröffentlichte Arbeit wird zugeschrieben. Der Zielwechsel zur
verallgemeinerten A381054-Aussage ist ausdrücklich; A383115 wird hier weder als
eigene neue Lösung noch als durch diesen Gegenbeweis entschieden verbucht.

\end{document}
